\section{The signed finite-group comparison}\label{sec:sr-finite}
This section gives the finite comparison underlying Anton's shorter route.
Every estimate in this section is proved; the existence of data satisfying
them on late intervals is stated separately in
Section~\ref{sec:sr-construction}.

Let \(H\) be a finite abelian group, \(n=|H|\), and put
\[
 \Gamma=(\mathbb Z/2)^2\times H,\quad
 \Phi=\{(0,b,h):b\in\mathbb Z/2,\ h\in H\},\quad f=|\Phi|=2n.
\]
Use additive notation for \(\Gamma\). Define
\(c=(1,1,0)\), \(t^\circ=(0,0,t)\),
\(\chi(a,b,h)=(-1)^a\), and \(\omega(a,b,h)=(-1)^b\).
The dual group \(\widehat\Gamma\) consists of its complex unitary characters;
its identity character is denoted \(1\). It has \(4n=2f\) elements.

For each rational prime \(p\), choose a subgroup \(I_p\subset\Gamma\) and
a Frobenius representative \(F_p\in\Gamma\).
Fix distinct odd primes \(q,d\). The host conditions are:
\[
 I_p=0\ (p\notin\{2,q,d\}),\qquad
 I_p\subset\Phi,\ F_p\in\Phi\ (p=2,q),\qquad
 I_d=\{0,(1,0,0)\}.
\]
For \(p^r\), \(r\geq1\), set
\[
 c_\psi(p^r)=\frac1{|I_p|}\sum_{h\in I_p}\psi(rF_p+h).
\]
Extend these coefficients multiplicatively to all positive integers.
Their absolute values are at most one, and \(c_1(m)=1\).
Choose arbitrary constants \(\ell_\psi\in\mathbb C\) for \(\psi\ne1\).
For \(y>0\), define
\[
 G_1(y)=E(y),\qquad
 G_\psi(y)=\sum_{r\leq y}\Lambda(r)c_\psi(r)(y/r-1)+y\ell_\psi
 \quad(\psi\ne1),
\]
using the same formula defining \(E\) also when \(0<y<1\).
Put
\[
 {\cal C}(y)=\Re(G_\chi(y)+G_{\chi\omega}(y)-G_\omega(y)),
 \qquad {\cal S}(y)=G_1(y)+G_\omega(y)-G_\chi(y)-G_{\chi\omega}(y).
\]
Thus \(\Re{\cal S}(y)=E(y)-{\cal C}(y)\).

\subsection{Orthogonality, including the ramified coefficients}
For every \(g\in\Gamma\), orthogonality on the finite abelian group gives
\begin{equation}\label{eq:sr-coset}
 \frac1f\sum_{\psi\in\widehat\Gamma}\psi(-g)c_\psi(p^r)
 =\begin{cases}2/|I_p|&g-rF_p\in I_p,\\0&\text{otherwise}.\end{cases}
\end{equation}
To verify this without any representation theory, interchange the
two finite sums. For each \(h\in I_p\), the inner sum is
\(\sum_\psi\psi(rF_p+h-g)\), equal to \(|\Gamma|\) if its argument is zero
and zero otherwise. The latter assertion follows by multiplying the sum
by a character nontrivial on the argument; such a character exists by
the decomposition of a finite abelian group into cyclic groups.
There is at most one \(h\) giving zero. Division by \(f|I_p|\) proves
\eqref{eq:sr-coset}.

In particular the raw four-character prime-power coefficient is
\[
 1+c_\omega(p^r)-c_\chi(p^r)-c_{\chi\omega}(p^r)
 =\frac4{|I_p|}
   \#\{h\in I_p:\chi(rF_p+h)=-1,\ \omega(rF_p+h)=1\}.
\]
This follows pointwise from
\(1+\omega-\chi-\chi\omega=(1-\chi)(1+\omega)\).
It is real and lies in \([0,4]\), even at ramified primes.
This positivity statement concerns the coefficients; it does not by itself
bound the complete signal, which also includes the normalizing terms in
the functions \(G_\psi\).

\subsection{The exact colour balance}
For \(t\in H\) write \(\rho_t(g)=c+g+t^\circ\).
It is a bijection from \(\Gamma\setminus\Phi\) to \(\Phi\).
Let \(a>0\), let \(w\) be positive on \(\Phi\), and choose an integer
\(N\geq a/w(\rho)\) for every \(\rho\in\Phi\). Define
\[
 \begin{split}
 {\cal A}(a,t,w)={}&
 2\sum_{\substack{p\leq N\ {\rm prime}\\p\notin\{2,q,d\}\\F_p\notin\Phi}}
  \log p
  \sum_{\substack{1\leq r\leq N\\rF_p\notin\Phi}}
       (a/p^r-w(\rho_t(rF_p)))_+\\
 &+\log d\sum_{r=1}^N
   \sum_{\substack{g-rF_d\in I_d\\g\notin\Phi}}
       (a/d^r-w(\rho_t(g)))_+ ,
 \end{split}
\]
and
\[
 M(a,w)=a\left(\log a-1-\gamma+\Re\ell_\chi
                    -\frac1f\sum_{\rho\in\Phi}\log w(\rho)\right).
\]
Finally put
\[
 B(a,t,w)=\Re\left[
 \frac1f\sum_{g\notin\Phi}\sum_{\psi\in\widehat\Gamma}
    \psi(-g)w(\rho_t(g))G_\psi(a/w(\rho_t(g)))\right].
\]
Then
\begin{equation}\label{eq:sr-balance}
 {\cal A}(a,t,w)-M(a,w)=B(a,t,w).
\end{equation}
Here is a direct verification. First expand the ramp sums in \(B\) and
use \eqref{eq:sr-coset}. Terms beyond \(N\) vanish by the choice of \(N\).
The terms at \(2,q\) vanish, because their Frobenius cosets lie in \(\Phi\).
For every other unramified prime the coefficient is two at the single
point \(rF_p\). At \(d\), its two-point coset has exactly one point outside
\(\Phi\), with coefficient one. These are precisely the two sums in
\({\cal A}\). No prime with \(p>N\), or exponent \(r>N\), contributes,
because \(p^r>N\) then and every relevant \(a/w\leq N\).

For the remaining terms, sum each character on the complement of
the index-two subgroup \(\Phi\):
\[
 \frac1f\sum_{g\notin\Phi}\psi(-g)
 =\begin{cases}1&\psi=1,\\-1&\psi=\chi,\\0&\text{otherwise}.\end{cases}
\]
This follows from the same orthogonality, first on \(\Phi\) and then
on its coset. The nontrivial-character linear terms therefore contribute
\(-a\ell_\chi\).
The normalizing terms in \(G_1\) contribute
\(-a\log a+(1+\gamma)a+(a/f)\sum_{\rho\in\Phi}\log w(\rho)\).
Their real parts give \(-M\), proving \eqref{eq:sr-balance}.

\subsection{Two uniform Lipschitz estimates}
For all \(A,w,w'>0\) and every character,
\begin{equation}\label{eq:sr-weightlip}
 |wG_\psi(A/w)-w'G_\psi(A/w')|
 \leq\frac{3A}{\min(w,w')}|w-w'|.
\end{equation}
For \(\psi\ne1\), write
\[
 wG_\psi(A/w)=\sum_r\Lambda(r)c_\psi(r)(A/r-w)_++A\ell_\psi .
\]
The linear term cancels. Each positive-part function is Lipschitz with
constant one, and only \(r\leq A/\min(w,w')\) can contribute.
Equation \eqref{eq:sr-cheb} proves \eqref{eq:sr-weightlip}.
For \(\psi=1\), the expression is the real continuous piecewise smooth
function
\[
 \sum_r\Lambda(r)(A/r-w)_+-A\log A+A\log w+(1+\gamma)A.
\]
Off its finitely many breakpoints in any compact interval of \(w>0\),
its derivative is \(A/w-\psi(A/w)\).
By \eqref{eq:sr-cheb}, its absolute value is at most \(2A/w\).
Integrating across the finitely many pieces proves the stronger bound
with constant two, and hence \eqref{eq:sr-weightlip}.

Also,
\begin{equation}\label{eq:sr-Elip}
 |E(x)-E(y)|\leq4|x-y|\qquad(x,y\geq1).
\end{equation}
For an integer \(m\geq1\), set \(L(m)=\sum_{r\leq m}\Lambda(r)/r\).
The factorial identity gives
\[
 \log(m!)=mL(m)-\sum_{r\leq m}\Lambda(r)\{m/r\},
\]
where the second sum lies between zero and \(\psi(m)\leq3m\).
The factorial integral bounds above imply
\(-1\leq L(m)-\log m\leq3\), using
\(1+\log m\leq m\).
If \(m\leq x<m+1\), then \(L(x)=L(m)\) and
\(0\leq\log x-\log m\leq\log2\). Therefore
\[
 -1-\log2\leq L(x)-\log x\leq3.
\]
The elementary harmonic integral inequalities give \(0\leq\gamma\leq1\).
Away from its integer breakpoints,
\(E'(x)=L(x)-\log x+\gamma\), whose absolute value is at most four.
Continuity and integration over the finitely many pieces between \(x,y\)
prove \eqref{eq:sr-Elip}.

\subsection{Transporting one signed bound}
Suppose \(x\geq3\), \(|y-x|\leq1\).
Apply \eqref{eq:sr-weightlip} with \(A=xy,w=x,w'=y\).
Taking real parts gives, for every \(\psi\),
\[
 |x\Re G_\psi(y)-y\Re G_\psi(x)|
 \leq3\max(x,y)|x-y|\leq4x.
\]
Using the upper inequalities for \(\chi,\chi\omega\) and the lower
one for \(\omega\), we obtain
\(x{\cal C}(y)\leq y{\cal C}(x)+12x\).
If \({\cal C}(x)\leq Bx^\theta\), \(B\geq0\), \(\theta\leq1\), then
\[
 {\cal C}(y)\leq Bx^\theta+B+12,
\]
because \(0\leq y\leq x+1\) and \(x^{\theta-1}\leq1\).
Together with \eqref{eq:sr-Elip} this proves
\begin{equation}\label{eq:sr-signaltransport}
 \Re{\cal S}(y)\geq E(x)-Bx^\theta-(B+16).
\end{equation}
The additive loss here uses the elementary constant four proved above.

\subsection{Averaging the thick part and bounding the repair}
Choose subsets \(T,T_2\subset H\), an integer \(m\geq1\),
numbers \(a_t>0\), a scale \(b>0\), and weights
\(w_{u,j}(\rho)>0\), \(u\in H,1\leq j\leq m,\rho\in\Phi\).
Call a point \(\rho\) thick for \(u\) if
\(\rho=(0,1,h+u)\) for some \(h\in T_2\).
Replace its weight by \(b\), keeping all other weights unchanged,
to obtain \(\widetilde w_{u,j}\).
Write \(\delta_2=|T_2|/n\). Define the thin mean
\[
 {\cal R}=\frac1n\sum_{u\in H}\sum_{t\in T}\sum_{j=1}^m
 \Re\left[\frac1f
 \sum_{\substack{g\notin\Phi\\\rho_t(g)\ {\rm not\ thick\ for}\ u}}
 \sum_{\psi\in\widehat\Gamma}
 \psi(-g)w_{u,j}(\rho_t(g))
 G_\psi(a_t/w_{u,j}(\rho_t(g)))\right].
\]
Then
\begin{equation}\label{eq:sr-thick}
 \frac1n\sum_{u,t,j}B(a_t,t,\widetilde w_{u,j})
 =\frac{\delta_2b}{2}\sum_{t\in T}\sum_{j=1}^m
       \Re{\cal S}(a_t/b)+{\cal R}.
\end{equation}
Indeed thickness requires that \(g=(1,0,h)\).
Averaging the translates \(T_2+u\) gives multiplicity \(|T_2|\)
at each \(h\in H\). Summing over \(h\) kills every character nontrivial
on \(H\). The four survivors \(1,\omega,\chi,\chi\omega\) have values
\(1,1,-1,-1\), respectively, at \((1,0,0)\).
Division by \(nf=2n^2\) leaves \(\delta_2/2\).
The complementary part is unchanged by the replacement, proving the identity.

Suppose moreover that \(|a_t/b-x|\leq1\), and the original thick weights
satisfy \(b\leq w_{u,j}(\rho)\leq b+\eta+1\), where \(\eta\geq0\).
For each fixed \(u,t,j\), \eqref{eq:sr-weightlip} bounds the change in
each weighted \(G_\psi\) by
\(3(x+1)(\eta+1)\).
There are \(f\) possible \(g\), \(2f\) characters, and a factor \(1/f\).
The triangle inequality therefore gives
\[
 |B(a_t,t,\widetilde w_{u,j})-B(a_t,t,w_{u,j})|
 \leq6f(x+1)(\eta+1).
\]
Consequently, if \(Q=|T|mb\) and
\begin{equation}\label{eq:sr-repair}
 6f(x+1)(\eta+1)\leq b/x,
\end{equation}
the total change in the averaged balance is at most \(Q/x\).

The thin mean can also be recovered from the original stronger
pointwise input. If at every thin node and for every \(\psi\),
\[
 |w_{u,j}(\rho)G_\psi(a_t/w_{u,j}(\rho))|\leq L,\qquad L\geq0,
\]
then the same count of terms gives
\begin{equation}\label{eq:sr-thin}
 {\cal R}\geq-2f|T|mL.
\end{equation}
If, more specifically, \(a_t\leq(x+1)b\),
\(a_t/w_{u,j}(\rho)\geq X'/3\) at thin nodes and
\(|G_\psi(z)|\leq z\vartheta\) for \(z\geq X'/3\), take
\(L=(x+1)b\vartheta\).
The inequality \(2f(x+1)\vartheta\leq1/x\) then yields
\({\cal R}\geq-Q/x\). No averaging cancellation is assumed in this argument.

\section{The exact geometric construction requirement}
\label{sec:sr-construction}
We now specify what must be supplied by the arithmetic geometry.
The distinction between a local implication and the existence statement
in this section is essential.

\begin{srdefinition}[A signed datum at \(x\)]
Fix \(C\geq0\) and \(0<\epsilon\leq1/2\). A signed datum consists of the
host and node data above, a cutoff \(N\), numbers \(\eta\geq0,k>0,A_E\geq0\)
and \({\rm Err}\), satisfying the following conditions.
Write \(t_0=|T|,t_2=|T_2|,Q=t_0mb\).
The host conditions hold, \(x\geq3\),
\(|a_t/b-x|\leq1\), and every \(a_t/w_{u,j}(\rho)\leq N\).
Thick weights lie in \([b,b+\eta+1]\), all weights are positive, and
\eqref{eq:sr-repair} holds. The three substantive inequalities are
\begin{align}
 \frac1n\sum_{u,t,j}\bigl({\cal A}(a_t,t,w_{u,j})-M(a_t,w_{u,j})\bigr)
       &\leq{\rm Err},\label{eq:sr-top}\\
 {\cal C}(x)&\leq3x^{1/2+\epsilon},\label{eq:sr-correction}\\
 {\cal R}&\geq-Q/x.\label{eq:sr-rest}
\end{align}
Finally the numerical shape requirements are
\begin{gather}
 A_E\leq2xb,\quad
 t_0\geq n x^{-\epsilon/8}/2,\quad
 x^{-\epsilon}\leq t_2/n\leq1,\label{eq:sr-shape1}\\
 x^{1/2-\epsilon}\leq m\leq x^{1/2+\epsilon/2},\quad
 xm\leq k\leq n^2x^{1/2-\epsilon},\quad
 26\leq x^{3\epsilon/8},\label{eq:sr-shape2}\\
 {\rm Err}\leq x^{C\epsilon}\left[
 t_0A_E(1+m^2/k)+m\left(2nmbx^\epsilon+
                 \frac{t_2kb x^\epsilon}{n^2x^{-\epsilon}}\right)\right].
 \label{eq:sr-error}
\end{gather}
\end{srdefinition}

\begin{srproposition}[Pointwise consequence]\label{prop:sr-point}
A signed datum at \(x\) gives \(E(x)\leq x^{1/2+(C+3)\epsilon}\).
\end{srproposition}
\begin{proof}
The condition \(26\leq x^{3\epsilon/8}\) implies
\(19\leq x^{1/2+\epsilon}\).
Apply \eqref{eq:sr-signaltransport} with \(B=3,\theta=1/2+\epsilon\)
and \(y=a_t/b\). It gives
\(\Re{\cal S}(a_t/b)\geq E(x)-4x^{1/2+\epsilon}\).
Equation \eqref{eq:sr-thick}, the thin lower bound, the repair bound,
the exact balance, and \eqref{eq:sr-top} therefore give
\begin{equation}\label{eq:sr-sizecomparison}
 Q\frac{\delta_2}{2}(E(x)-4x^{1/2+\epsilon})
       \leq{\rm Err}+2Q/x.
\end{equation}
All quantities divided by below are positive: \(n,b,m,t_0,k,\delta_2\)
are positive by the shape requirements.

Put \(Y=x^{1/2+13\epsilon/8}\) and \(P_0=x^{C\epsilon}\).
Divide the bracket in \eqref{eq:sr-error} by \(Q\).
Its four terms satisfy, respectively,
\[
 \frac{A_E}{mb}\leq2x^{1/2+\epsilon},\qquad
 \frac{A_Em}{bk}\leq2,\qquad
 \frac{2nmx^\epsilon}{t_0}\leq4x^{1/2+13\epsilon/8},\qquad
 \frac{t_2kx^{2\epsilon}}{n^2t_0}\leq2x^{1/2+9\epsilon/8}.
\]
These follow in order from \(A_E\leq2xb\) and the lower bound on \(m\);
from \(k\geq xm\); from the lower bound on \(t_0\) and upper bound on \(m\);
and from the upper bounds on \(t_2,k\) and lower bound on \(t_0\).
Each right side is at most its numerical coefficient times \(Y\), so
\({\rm Err}\leq10P_0YQ\).
Dividing \eqref{eq:sr-sizecomparison} by \(Q\delta_2/2\) gives
\[
 E(x)\leq4x^{1/2+\epsilon}
       +20P_0Yx^\epsilon+4x^{\epsilon-1}
 \leq4x^{1/2+\epsilon}+22P_0Yx^\epsilon .
\]
For the last step use \(x\geq3\) and \(P_0Y\geq1\).
Since \(C\geq0\),
\(x^{1/2+\epsilon}\leq P_0Yx^\epsilon
=x^{1/2+(C+21/8)\epsilon}\).
The result is at most \(26x^{1/2+(C+21/8)\epsilon}\).
Absorb 26 by \(x^{3\epsilon/8}\), obtaining the asserted bound.
\end{proof}

\begin{srdefinition}[Uniform supply]\label{def:sr-supply}
The supply assertion is the following precise statement:
there exist \(C\geq0,\epsilon_0>0\), with \(\epsilon_0\leq1/2\), such that
for every \(0<\epsilon<\epsilon_0\) and every \(Y_0\), there is
\(Y\geq\max(Y_0,1)\) for which a signed datum with the same \(C,\epsilon\)
exists at every fourth power \(j^4\in[Y,Y^{3/2}]\).
The host, group, field, weights and auxiliary parameters may vary with
the sample. The constant \(C\) and the order of quantifiers may not.
\end{srdefinition}

\begin{srtheorem}[Signed supply implies RH]\label{thm:sr-supply}
The assertion of Definition~\ref{def:sr-supply} implies ordinary RH.
\end{srtheorem}
\begin{proof}
Given any \(\alpha>0\), choose \(0<\epsilon<\epsilon_0\) with
\((C+3)\epsilon\leq\alpha\).
Proposition~\ref{prop:sr-point} bounds \(E(j^4)\) by
\((j^4)^{1/2+\alpha}\) at all the supplied samples.
By \eqref{eq:sr-integer} and \eqref{eq:sr-discard},
\[
 P(j^4)\geq\gamma j^4-(1+K)(j^4)^{1/2+\alpha}.
\]
The constant \(1+K\) is uniform across the chosen interval.
Theorem~\ref{thm:sr-prime} now applies.
\end{proof}

\paragraph{Status of the construction.}
The preceding theorem proves an implication. This edition does not prove
Definition~\ref{def:sr-supply}. In particular it does not construct
the geometric witnesses giving \eqref{eq:sr-top} with the error
\eqref{eq:sr-error} on arbitrarily late intervals. It also does not give
an independent proof that the signed correction bound
\eqref{eq:sr-correction} and the complementary bound
\eqref{eq:sr-rest} can be achieved simultaneously with those witnesses.
The archive-to-signed-data map transfers these requirements from stronger
archive inputs; it does not remove the requirement to prove their existence.
Consequently Theorem~\ref{thm:sr-supply} cannot be described as an
unconditional proof supplied by this paper.
